\chapter{Keeping the Lean verifier from drifting away from the deployed one}\label{ch:drift}

The requirement, stated by the owner: a verifier that omits a check leanVM makes, or makes a
weaker one, must have no proof of knowledge soundness; a verifier that makes a wrong check, or
one leanVM does not make, must have no proof of completeness. The theorems, not the tests, are
to keep leanerVM faithful. This chapter says where that holds today, where it does not and why,
and what would make it hold as the work proceeds and the pin moves.

\section{Where a weakening can hide}\label{sec:drift-hiding}

The chain from \lean{verify} to the semantics (\cref{fig:chain}) has five places where a
missing or wrong check could go unnoticed by the master theorems. For each, the review says
whether it occurs today.

\paragraph{The relation weakened in step with the verifier.} The relation is
\lean{M3Holds I input q}, and \lean{I} is data: its \lean{publicLines}, its \lean{aux}, its
\lean{counts}, its \lean{boundary}, its \lean{constraints} and \lean{flushes}, its
\lean{layout}. A verifier that omits the claim on the top memory limb, together with an
instance that lists two public lines instead of three, is knowledge sound: every theorem of the
spine and of the public-input phase holds of the two-line instance. A Flock verifier that checks
nothing, together with an instance whose \lean{aux} is \lean{True}, is knowledge sound at error
zero (probe \code{NoCheckFlock}). An instance whose layout aliases columns has an empty
relation and both master theorems (probe \code{P2Relation}). \emph{This occurs.} It is the
price of stating the protocol over an abstract instance: the instance is where leanVM's
relation enters, and the oracle protocol's theorems cannot check the instance. What catches a
weakened instance is the adaptor's soundness theorem, \lean{satisfiedBy_witnessOf}, whose
conclusion \lean{SatisfiedBy} does not weaken with the instance: a two-line instance leaves its
public-word conjunct unprovable (\lean{SatisfiedBy.word0\_eq} has the top limb zero by the type
of \lean{PublicInput}); a trivial \lean{aux} leaves \lean{Blake2sRowsValid} unprovable; an
aliasing layout leaves the completeness theorem \lean{m3Holds_stackOf} unprovable. Behind the
adaptor, \lean{constraintSoundness} anchors \lean{SatisfiedBy} to \lean{ValidExecution}. None of
these three theorems is built. \textbf{Until the adaptor is built, the master theorems
constrain the instance by nothing, and a weakening of the instance is caught by no theorem on
\code{main}.}

\paragraph{A property supplied as a hypothesis that nothing discharges.} The master theorems
take \lean{Phases.Complete} and \lean{Phases.Security} as hypotheses; four of the five phases
are unbuilt, so the hypotheses are undischarged. Worse, \lean{Phases.Security} can be
discharged by phases that check nothing at declared error one (probe \code{P4Junk}): the error
is the phases' own field. \emph{This occurs}, and it is the finding that most changes how the
soundness theorem should be read (\cref{sec:nv-hypotheses}). The repair is to fix the error in
the spine, so that a phase can only meet its slot by earning the bound.

\paragraph{A check behind an assumed interface.} Flock, ring switching, WHIR's soundness, the
Merkle compilation and Fiat--Shamir are consumed as interfaces. Today none exists in Lean; when
they exist, a check the deployed verifier makes inside one of them (the lincheck's terminal
identity; the Merkle path check; the grinding check) is load-bearing only for the interface's
own theorem, which another roadmap or an upstream library owes. \emph{This occurs by design},
and the report accepts it on two conditions: that each interface is a structure whose fields are
theorems with named owners and no inhabitant in the repository until they are proved (the
repository's own rule against uninhabited load-bearing types then flags them), and that the
interfaces are stated for leanVM's constructions, which the Fiat--Shamir and BCS ones today are
not (\cref{sec:faith-opening}).

\paragraph{A property lost between the abstract instance and leanISA's, or in the adaptor.}
The instance for leanISA does not exist; the adaptor does not exist. What the review found is
that the blueprint's sketch of them cannot be written as it stands (the admissibility predicate,
the strided reader of the limb columns, the Flock predicate, the fit of the layout) and that
the sketch's \lean{witnessOf} builds parts of the witness from the program, which is sound only
if balance is then proved of the mixed witness (\cref{ch:faithfulness}). \emph{Cannot occur
yet}, since nothing is there to lose it; the report's recommendations for Layer 3 are where it
must be prevented.

\paragraph{A check present in the oracle verifier and absent from the compiled one, or in Lean
and absent from Rust.} The compiled verifier does not exist. Between Lean and Rust the review
found one check Lean makes in a stricter form than the deployed verifiers (the per-limb
public-input check, where they check one combined equation), one that the deployed verifiers make and Lean's oracle protocol cannot (the proof of
work), and three format checks the blueprint does not list (canonical encodings of the sizes
and of the root, the stream fully consumed). \emph{This occurs} at the public-input phase
today, and would occur at the compilation as planned.

\section{The public-input phase, as the test case}\label{sec:drift-pubinput}

The owner's own example was the top lane of the two public words. Following it to its anchor:

\begin{enumerate}
\item In the deployed verifier, two things enforce it: the statement-level rejection of a
  public word with a nonzero top limb at setup (\file{cpu/mod.rs:141-143}), and the pooled claim
  that the top memory limb is zero at $(r, 0, \dots, 0)$, for which no scalar is sent.
\item In Lean, the first is enforced by the type: leanISA's \lean{PublicInput} is four lanes of
  $\K$, so a nonzero top limb is not writable. The second is the third public line of the
  leanISA instance, which the blueprint specifies and which nothing on \code{main} states.
\item The phase pools the third claim because the instance lists a third line. Drop the line
  and every theorem of the spine and of the phase still holds (\cref{sec:nv-pubinput}). The
  master theorems do not catch it.
\item The adaptor's theorem \lean{satisfiedBy_witnessOf} would: \lean{SatisfiedBy} requires
  the memory cells $0$ and $1$, as three-limb values, to be the public words, whose top limb is
  zero by type. Without the third line the top limb of the extracted memory is unconstrained
  and the conjunct is unprovable. Behind it, \lean{constraintSoundness} carries the conjunct to
  the semantics.
\item So the answer to the owner's question is: \textbf{the master theorems do not catch it;
  the adaptor's soundness theorem does, once it exists; the blueprint's acceptance test 10 is a
  test on the pool, not a theorem.} The chain is anchored, but its anchor is two unbuilt
  theorems away from the oracle protocol.
\end{enumerate}

The same phase gave a second, harder lesson. Its check on the prover's message is not
load-bearing for any theorem, because the verifier pools the values it computes rather than the
values it received (\cref{sec:nv-pubinput}). The requirement cannot be met for that check by
any theorem about the phase as designed, since omitting the check yields a sound verifier: a
verifier that ignores the message and pools the statement's values is right to do so. What
makes the check load-bearing is the design of the deployed verifiers, which pool the values
\emph{sent}: with \lean{pooledFrom} in place of \lean{pooled} the omission and the weakening are
refuted by Lean theorems and a wrong check breaks completeness for every prover (probes
\code{Probe2TrustProver}, \code{Probe3b}, \code{Probe7a}). The rule this gives is general:
\emph{a check is load-bearing only if the verifier commits to what it checked.} A phase that
recomputes what the prover sent makes the message, and its check, decorative.

\section{The checks of the deployed verifier, one by one}\label{sec:drift-checks}

For every check the deployed verifier makes at the pin, the tables below give where it is in
the three sources, what it protects, the concrete cheating prover or witness a verifier
without it would accept, and, where a phase is built, the Lean check and the theorem that
makes it necessary. They are the review's enumeration from the specification and from the
rejection paths of the Rust and Python verifiers; the setup's checks and the format checks are
included, since the compiled verifier must make them and the blueprint does not list them.

\InputIfFileExists{sections/gen/checks-bus}{}{\emph{[The checks of the setup, the commitment
and the bus phase: fragment pending.]}}
\InputIfFileExists{sections/gen/checks-table-pub}{}{\emph{[The checks of the table sumcheck
and the public-input phase: fragment pending.]}}
\InputIfFileExists{sections/gen/checks-flock-ring}{}{\emph{[The checks of the BLAKE2s validity
phase and ring switching: fragment pending.]}}
\ifopeningready\InputIfFileExists{sections/gen/checks-opening}{}{}\else\emph{[The checks of the opening and of the compiled verifier: fragment pending.]}\fi

Three checks stand out because no planned theorem makes them load-bearing: the proof of work
(grinding), the canonical encodings, and the consumption of the stream. The first is a soundness
check in the non-interactive setting only, and no published theorem covers Fiat--Shamir with
proof of work for a multi-round protocol; the other two are format checks whose omission lets
\lean{verify} accept streams the deployed verifier rejects. All three need a mutation test in
the compiled verifier's fixture (\cref{sec:drift-mechanisms}).

\section{What keeps it true}\label{sec:drift-mechanisms}

\paragraph{Stated once, so that it cannot drift.} The relation, the seams and the error: the
first two are the spine's already; the error is to be (\cref{sec:aud-proposals}). With the
error fixed by the spine, a phase that declares less than it earns does not typecheck against
its slot.

\paragraph{What each phase delivers before it lands.} For each check of the phase's verifier,
a \emph{refutation}: a Lean theorem that the verifier without that check has no round-by-round
knowledge soundness against the seams at any error below one, whatever the extractor and the
state function; and for each check the deployed verifier does not make, a theorem that the
verifier with it is not perfectly complete. The public-input phase's audit gives the two
generic lemmas this needs (\code{not\_rbr} and \code{not\_perfectCompleteness},
\cref{ch:probes}); they belong in \file{ToArkLib/}. A mutation whose omission cannot be
refuted is a check the verifier does not rely on, and the phase's design should be revised
until it does (the rule of \cref{sec:drift-pubinput}). These refutations are tests under
\file{tests/}, run by \code{lake test}, and a phase's pull request lists them.

\paragraph{What ties the Lean verifier to the deployed one.} The compiled verifier's refinement
theorem ties \lean{verify} to the oracle protocol by definition; only a fixture ties either to
the Rust. The fixture is a proof produced by the pinned Rust prover, accepted by \lean{verify}
under \lean{\#guard}, with a mutation per check: each of the checks of \cref{sec:drift-checks}
gets a mutated proof that the deployed verifier rejects and that \lean{verify} must reject too.
The blueprint's six mutations become as many as there are checks. Reverse tests, transcripts
the specification rejects run through the Rust and Python verifiers, are how the per-limb
public-input divergence would have been found by a test rather than by a reader; the review
recommends them for every check whose statement in the specification differs from the code.

\paragraph{What is examined again when the pin moves.} The leanVM pin: every check of
\cref{sec:drift-checks} against the new verifier's rejection paths, the message and challenge
counts of every phase, the layouts and their tie-breaking orders, the constants of Flock and
of Fiat--Shamir, the parameter tables. The library pins: the admitted theorems the ledger names
(this upgrade lifted none), the notation of the probability bridges (this upgrade changed it),
the meaning of numerals in $\K$ (this upgrade changed it), and every acceptance test that
\lean{decide}s a numeral.

\paragraph{What the repository's validation can enforce.} Today it enforces the axiom
closure, the absence of \lean{sorry}, the layer DAG and the documentation index. It can also
enforce: the wall (a rule in \file{scripts/check-layers.sh} listing the modules allowed to
import the arithmetization, which today would flag \file{Protocol/Basic.lean}); the presence,
for every built phase, of a refutation test per check (a script that reads the phase's checks
from a manifest and asserts a test per entry); the fixture and its mutations, once the
compiled verifier exists; and, at a pin move, a script that re-derives the message and
challenge counts of every phase from its \lean{Def} and compares them with counts dumped from
the Rust prover. What it cannot enforce is that a phase's definition is leanVM's: that rests
on review against the specification, which is why this report's tables of the transcript
exist.
